Este trabajo final de máster aborda el desarrollo y la ampliación del marco de 
simulación dinámica de sistemas de potencia dentro de la plataforma VeraGrid, el software de código abierto de eRoots,
 con el objetivo de incorporar 
tanto el análisis de pequeña señal como los fundamentos de la simulación de transitorios electromagnéticos (EMT). 
Todo el desarrollo se ha llevado a cabo en lenguaje Python, poniendo énfasis en la modularidad y escalabilidad del 
código mediante una arquitectura clara y orientada a la evolución futura de la herramienta. El trabajo se enmarca en el 
contexto actual de los sistemas eléctricos, caracterizados por una creciente presencia de dispositivos basados en electrónica 
de potencia y por la necesidad de herramientas de simulación capaces de representar fenómenos dinámicos en diferentes escalas 
temporales.

En una primera parte, se desarrolla un módulo de análisis de pequeña señal para modelos dinámicos en dominio de valor de tensión eficaz (RMS, por sus siglas en inglés, root mean square), basado 
en la formulación simbólica de ecuaciones diferenciales-algebraicas y en su linealización alrededor del punto de operación. 
Esta aproximación permite el cálculo automático de jacobianos, autovalores y factores de participación, garantizando la 
coherencia entre la simulación no lineal y el análisis linealizado del sistema, y facilitando el estudio de la estabilidad y 
del comportamiento modal del sistema eléctrico. Los resultados obtenidos se validan mediante la comparación con ANDES, 
una herramienta de referencia en el ámbito simbólico-numérico.

Posteriormente, se establecen los fundamentos para la simulación EMT mediante la implementación de modelos en el dominio 
temporal con variables instantáneas en coordenadas abc. Se desarrollan modelos básicos de generador síncrono, 
líneas de transmisión y cargas, con especial énfasis en el modelo de línea basado en el modelo de Bergeron y en la 
teoría de ondas viajeras. La formulación EMT se implementa utilizando el algoritmo de Dommel conjuntamente con el marco 
de trabajo simbólico-numérico y esquemas de integración trapezoidal, asegurando una representación físicamente coherente 
de los fenómenos electromagnéticos rápidos. La validación se realiza mediante diferentes casos de estudio, que permiten 
analizar la propagación de ondas y los mecanismos de reflexión bajo distintas condiciones.

Los resultados obtenidos muestran un comportamiento basado en la teoría clásica de los transitorios electromagnéticos y 
confirman la corrección tanto de la formulación matemática como de la implementación numérica. 
Este trabajo establece una base sólida, extensible y robusta para el desarrollo de un motor de simulación EMT 
avanzado integrado dentro de VeraGrid, con un alto potencial de impacto tanto en el ámbito académico como en aplicaciones 
industriales reales.
